\Askhsh{18565_SOLUTION}{1}
\begin{ylh}{1}
    \item [Γεωμετρία, Πυθαγόρειο Θεώρημα, Εμβαδά]
\end{ylh}
\λύση
\begin{wrapfigure}[24]{r}{7.5cm} % Adjust lines and width as needed
    \centering
    \begin{tikzpicture}[scale=0.7]
        % Define points for the common external tangent configuration
        % E at origin, O up, K right, A up from O, B up from K
        \tkzDefPoint(0,0){E}
        \tkzDefPoint(0,5){O} % OE = 5
        \tkzDefPoint(0,7){A} % AE = 2, OA = 7
        \tkzDefPoint(12,0){K} % KE = 12
        \tkzDefPoint(12,2){B} % KB = 2 (This creates rectangle ABKE, so AB=KE=12)

        % Draw segments OA, KB (radii)
        \tkzDrawSegment[l3](O,A)
        \tkzDrawSegment[l3](K,B)

        % Draw tangent AB
        \tkzDrawSegment[l3](A,B)

        % Draw construction line KE (dashed)
        \tkzDrawSegment[l4,dashed](K,E)

        % Draw line OK (distance between centers)
        \tkzDrawSegment[l3](O,K)

        % Mark right angles
        \tkzMarkRightAngle(O,A,B)
        \tkzMarkRightAngle(K,B,A)
        \tkzMarkRightAngle(O,E,K)

        % Labels for points
        \tkzLabelPoint[below left](O){O}
        \tkzLabelPoint[above left](A){A}
        \tkzLabelPoint[left](E){E}
        \tkzLabelPoint[below right](K){K}
        \tkzLabelPoint[above right](B){B}

        % Labels for lengths
        \tkzLabelSegment[left](O,E){5}
        \tkzLabelSegment[left](E,A){2}
        \tkzLabelSegment[right](K,B){2}
    \end{tikzpicture}
    % Separate visualization for the length breakdown of OK
    \begin{tikzpicture}[xscale=0.3] % Scale x to show relative lengths
        \tkzDefPoint(0,0){O_line}
        \tkzDefPoint(7,0){G_line}
        \tkzDefPoint(11,0){D_line}
        \tkzDefPoint(13,0){K_line}
        \tkzDrawSegment(O_line,K_line)
        \tkzDrawPoints(O_line,G_line,D_line,K_line)
        \tkzLabelPoint[below](O_line){O}
        \tkzLabelPoint[below](G_line){$\Gamma$}
        \tkzLabelPoint[below](D_line){$\Delta$}
        \tkzLabelPoint[below](K_line){K}
        \tkzLabelSegment[above](O_line,G_line){7}
        \tkzLabelSegment[above](G_line,D_line){4}
        \tkzLabelSegment[above](D_line,K_line){2}
    \end{tikzpicture}
\end{wrapfigure}
\textbf{α)}
\begin{rlist}
    \item Οι ακτίνες $OA$ και $KB$ είναι κάθετες στο εφαπτόμενο τμήμα $AB$ στα σημεία επαφής $A$ και $B$. Το τμήμα $KE$ είναι παράλληλο στο $AB$ και αφού $AB \perp OA$, θα είναι και $KE \perp OA$. Άρα το τετράπλευρο $ABKE$ έχει 3 ορθές γωνίες και είναι ορθογώνιο. Επομένως $AE=KB=2$ και $OE=OA-EA=7-2=5$. Στο ορθογώνιο τρίγωνο $OEK$ η υποτείνουσά του $OK$ είναι $O\Gamma+\Gamma\Delta+\Delta K=7+4+2=13$. Εφαρμόζοντας το Πυθαγόρειο Θεώρημα έχουμε: $KE^2=OK^2-OE^2$ ή $KE^2=13^2-5^2$ ή $KE^2=169-25$ ή $KE^2=144$, άρα $KE=12$. Άρα $AB=KE=12$.
    \item Στο τετράπλευρο $ABKO$ είναι $OA \perp AB$ και $KB \perp AB$, άρα $OA \parallel KB$. Επίσης $OA=7 \neq 2=KB$, άρα οι απέναντι πλευρές του είναι παράλληλες και άνισες, οπότε δεν είναι παραλληλόγραμμο. Άρα το $ABKO$ είναι τραπέζιο, στο οποίο η μικρή του βάση είναι η $KB$, η μεγάλη βάση του η $OA$ και το ύψος του είναι το μήκος του εφαπτόμενου τμήματος $AB$. Οπότε $(\text{ABKO}) = \frac{KB+OA}{2} \cdot AB = \frac{2+7}{2} \cdot 12 = 54 \text{ τ.μ.}$
\end{rlist}
\textbf{β)} Το $ABKE$ είναι ορθογώνιο, άρα $(\text{ABKE}) = AB \cdot KB$. Επειδή $(\text{ABKE}) = 4\sqrt{14}$, έχουμε:
$4\sqrt{14} = AB \cdot 2$ άρα $AB = 2\sqrt{14}$. Από το α) ερώτημα $AB=KE$, οπότε $KE=2\sqrt{14}$ και από το Π.Θ. στο τρίγωνο $OKE$ είναι: $OK^2 = OE^2+KE^2$ ή $OK^2=5^2+(2\sqrt{14})^2$ , ή $OK^2 = 25+56$, ή $OK^2=81$, άρα $OK=9$. Για το τμήμα $OK$ ισχύει ότι $OK=O\Gamma+\Gamma\Delta+\Delta K$, ή $9=7+\Gamma\Delta+2$, δηλαδή $9=9+\Gamma\Delta$, οπότε $\Gamma\Delta=0$.
Δηλαδή η διάκεντρος $OK$ των δύο κύκλων ισούται με το άθροισμα των ακτινών τους. Άρα οι κύκλοι εφάπτονται εξωτερικά.